\section{The exact affine Boussinesq wave and its viscous counterpart}
\label{sec:boussinesq-core}

This section reconstructs the local calculation in Alp\"oge and Buckmaster's
\emph{Blowup for the Boussinesq equations with smooth forcing}, Sections 1.2
and 3.1, from its equations. The source is
\url{https://cims.nyu.edu/~tristanb/boussinesq.pdf}.
The pressure construction and the diffusion calculation below make explicit
what the local wave does in the momentum equation. All coordinates, the
upward buoyancy direction, the frequency parameter, and the signed
amplitudes are retained. The fields in this section are defined on
$\mathbb R^2$; their lack of spatial decay is part of their stated scope.

\subsection{Coordinates, operators, and prescribed affine forces}

Write $x=(x_1,x_2)$, $e_2=(0,1)$, and
\[
J=\begin{pmatrix}0&-1\\1&0\end{pmatrix},\qquad
\nabla^\perp=J\nabla,\qquad
\operatorname{curl}(w_1,w_2)=\partial_1w_2-\partial_2w_1.
\]
Thus $\operatorname{curl}(\nabla^\perp\psi)=\Delta\psi$.
Let $I\subset\mathbb R$ be an open interval containing $0$.
Let $D:I\to\mathbb R^{2\times2}$ and $G:I\to\mathbb R^2$ be smooth,
with $\operatorname{tr}D=0$. Put
\[
u_b(x,t)=D(t)x,\qquad \theta_b(x,t)=G(t)\cdot x.
\]
For any prescribed smooth pressure $p_b:I\times\mathbb R^2\to\mathbb R$,
define, without discarding any term,
\begin{align}
f_{\theta,b}(x,t)&=(\dot G+D^TG)\cdot x,\label{eq:bc-base-scalar}\\
f_{u,b}(x,t)&=(\dot D+D^2)x+\nabla p_b-(G\cdot x)e_2.
\label{eq:bc-base-vector}
\end{align}
For constants $\kappa,\nu\geq0$ these fields satisfy
\begin{align}
\partial_t\theta_b+u_b\cdot\nabla\theta_b
 &=\kappa\Delta\theta_b+f_{\theta,b},\label{eq:bc-scalar-pde}\\
\partial_tu_b+(u_b\cdot\nabla)u_b+\nabla p_b
 &=\nu\Delta u_b+\theta_be_2+f_{u,b},\qquad
\operatorname{div}u_b=0. \label{eq:bc-momentum-pde}
\end{align}
Indeed, $\partial_tu_b=\dot D x$, $(u_b\cdot\nabla)u_b=D^2x$,
$\nabla\theta_b=G$, and both affine Laplacians vanish. These substitutions
give exactly \eqref{eq:bc-base-scalar}--\eqref{eq:bc-base-vector}.
In particular, an arbitrary time-dependent pair $(D,G)$ is not being
silently declared to be an unforced background.

The source uses the letter $\kappa$ for $|\zeta|^2$ inside its local lemma.
Here $q=|\zeta|^2$ denotes that exact quantity, while $\kappa$ denotes
the additional thermal diffusion coefficient. Neither quantity is set
equal to one.

\subsection{A general periodic profile and transported phase}

Fix $\lambda>0$ and $\zeta_0\in\mathbb R^2\setminus\{0\}$.
The phase vector is the solution of
\begin{equation}
\dot\zeta=-D^T\zeta,\qquad \zeta(0)=\zeta_0.
\label{eq:bc-phase}
\end{equation}
It stays nonzero on $I$. Here is a direct justification of the linear
ODE facts used throughout this section. On a compact time interval the
matrix coefficient is bounded by some finite $L$. Successive integration
of the equation $M'=DM$, $M(0)=\mathrm{Id}$ gives a matrix series whose
$n$th term has norm at most $(L|t|)^n/n!$. It and its differentiated
integral equation converge uniformly on that interval, producing a
solution. The same construction solves $N'=-ND$, $N(0)=\mathrm{Id}$.
Differentiation gives $(NM)'=0$, so $NM=\mathrm{Id}$. Smoothness follows
by repeated differentiation of the equation. For uniqueness, the
difference of two solutions is bounded, and iterating its homogeneous
integral equation bounds its norm by a fixed constant times
$(L|t|)^n/n!$ for every $n$; this tends to zero. The construction on
overlapping compact intervals consequently agrees. Then
$\zeta=M^{-T}\zeta_0$ gives \eqref{eq:bc-phase} and cannot vanish.
The same integral-series argument applies to the amplitude systems below.

Let $F:\mathbb R\to\mathbb R$ be smooth, odd, $2\pi$-periodic,
and not identically zero. Its integral over one period is zero:
the integral on $[-\pi,\pi]$ vanishes by oddness, and periodicity makes
all period integrals equal. Therefore
\[
P_0(s)=\int_0^sF(a)\,da,\qquad
P(s)=P_0(s)-\frac1{2\pi}\int_{-\pi}^{\pi}P_0(a)\,da
\]
defines a smooth, even, mean-zero, periodic primitive with $P'=F$.
For evenness, substitution $a=-b$ gives $P_0(-s)=P_0(s)$.
For periodicity,
$P_0(s+2\pi)-P_0(s)=\int_s^{s+2\pi}F(a)\,da=0$.
Subtracting the displayed constant supplies zero mean without changing
the derivative.

Set
\[
q(t)=|\zeta(t)|^2>0,\qquad
s(x,t)=\lambda\zeta(t)\cdot x,\qquad
C(t)=\frac{\Omega(t)}{\lambda q(t)}.
\]
The signed amplitudes are $\Theta,\Omega:I\to\mathbb R$.
Define the added fields
\begin{equation}
\vartheta=\Theta F(s),\qquad
\psi=\frac{\Omega}{\lambda^2q}P(s),\qquad
v=CJ\zeta F(s),\qquad
\varpi=\Omega F'(s).
\label{eq:bc-wave}
\end{equation}
Every identity below concerns these original fields.

\begin{proposition}[Exact uncut inviscid wave, including pressure]
\label{prop:bc-uncut}
For $\kappa=\nu=0$, take the unique smooth amplitude solution with any
real initial amplitudes and
\begin{equation}
\dot\Theta=-\frac{J\zeta\cdot G}{\lambda q}\Omega,\qquad
\dot\Omega=\lambda\zeta_1\Theta.
\label{eq:bc-amplitudes}
\end{equation}
Define
\begin{equation}
h=\dot C J\zeta+2C D J\zeta-\Theta e_2,\qquad
\pi(x,t)=-\frac{\zeta\cdot h}{\lambda q}P(s(x,t)).
\label{eq:bc-pressure}
\end{equation}
Then
\[
\theta=\theta_b+\vartheta,\qquad u=u_b+v,\qquad p=p_b+\pi
\]
satisfy the full inviscid scalar and momentum equations
\eqref{eq:bc-scalar-pde}--\eqref{eq:bc-momentum-pde} with exactly the
same prescribed forces \eqref{eq:bc-base-scalar}--\eqref{eq:bc-base-vector}.
Their vorticity is
$\omega=D_{21}-D_{12}+\varpi$.
\end{proposition}

\begin{proof}
First $\nabla s=\lambda\zeta$. Since $P'=F$,
\[
\nabla^\perp\psi=\frac{\Omega}{\lambda q}J\zeta F(s)=v,\qquad
\Delta\psi=\frac{\Omega}{\lambda^2q}\lambda^2q F'(s)=\varpi.
\]
Also
\[
\operatorname{div}v=C\lambda(J\zeta\cdot\zeta)F'(s)=0.
\]
Write $L_b=\partial_t+(Dx)\cdot\nabla$. The precise phase calculation is
\[
L_bs=\lambda\dot\zeta\cdot x+\lambda\zeta\cdot Dx
=\lambda(\dot\zeta+D^T\zeta)\cdot x=0.
\]
The wave gradients are
\[
\nabla\vartheta=\lambda\Theta\zeta F'(s),\qquad
\nabla\varpi=\lambda\Omega\zeta F''(s).
\]
Their scalar products with $v$ vanish because $J\zeta\cdot\zeta=0$.
The background vorticity $D_{21}-D_{12}$ is spatially constant, so
$v\cdot\nabla\omega_b=0$ as well. Consequently the scalar residual
increment is exactly
\[
L_b\vartheta+v\cdot\nabla\theta_b+v\cdot\nabla\vartheta
=\left(\dot\Theta+\frac{\Omega}{\lambda q}J\zeta\cdot G\right)F(s)=0.
\]
The vorticity residual increment is exactly
\[
L_b\varpi+v\cdot\nabla\omega_b+v\cdot\nabla\varpi-\partial_1\vartheta
=(\dot\Omega-\lambda\zeta_1\Theta)F'(s)=0.
\]
Thus far this is the scalar--vorticity calculation. We next prove its
explicit lift to the momentum equation.

Write
$D=\left(\begin{smallmatrix}a&b\\c&-a\end{smallmatrix}\right)$.
Multiplication gives $DJ+JD^T=0$. Hence
$J\dot\zeta=-JD^T\zeta=DJ\zeta$. The phase identity and
$v\cdot\nabla v=0$ give
\begin{align*}
L_bv+(v\cdot\nabla)u_b+(v\cdot\nabla)v-\vartheta e_2
&=(\dot C J\zeta+CJ\dot\zeta+CDJ\zeta-\Theta e_2)F(s)\\
&=hF(s).
\end{align*}
To show that the pressure in \eqref{eq:bc-pressure} cancels this exact
vector, compute
\[
\dot q=-2\zeta\cdot D^T\zeta=-2\zeta\cdot D\zeta,\qquad
(J\zeta)\cdot DJ\zeta=-\zeta\cdot D\zeta=\frac{\dot q}{2}.
\]
The second equality follows by multiplication with the displayed
trace-free matrix, without symmetry assumptions on $D$. Therefore
\[
(J\zeta)\cdot h
=\dot Cq+C\dot q-\Theta\zeta_1
=\frac{\dot\Omega}{\lambda}-\Theta\zeta_1=0.
\]
The vectors $\zeta,J\zeta$ are orthogonal and have the same positive
squared length $q$, so
$h=q^{-1}(\zeta\cdot h)\zeta$.
Finally
\[
\nabla\pi=-\frac{\zeta\cdot h}{q}\zeta F(s)=-hF(s).
\]
The momentum residual increment consequently vanishes component by
component. The already computed divergence and curl establish the
remaining claims.
\end{proof}

If $F(s)=\sin s$, the mean-zero primitive is $P(s)=-\cos s$.
The signs in \eqref{eq:bc-wave} then read
\[
\psi=-\frac{\Omega}{\lambda^2q}\cos s,\qquad
v=\frac{\Omega}{\lambda q}J\zeta\sin s,\qquad
\varpi=\Omega\cos s,
\]
so they agree with the sine-wave presentation, including its negative
streamfunction coefficient.

The profile can also be chosen exactly linear near zero. For completeness,
fix $0<a<b<\pi$, put $\rho(t)=e^{-1/t}$ for $t>0$ and $\rho(t)=0$
for $t\leq0$, and set
\[
\chi(s)=\frac{\rho(b^2-s^2)}
{\rho(b^2-s^2)+\rho(s^2-a^2)},\qquad
F(s)=\sum_{m\in\mathbb Z}(s-2\pi m)\chi(s-2\pi m).
\]
The denominator is strictly positive: its two arguments cannot both be
nonpositive because $a<b$. The function $\rho$ is smooth at zero
since each one-sided derivative is a polynomial in $1/t$ times
$e^{-1/t}$ and tends to zero as $t\downarrow0$.
Thus $\chi$ is smooth and even, equals one for $|s|\leq a$, and
vanishes for $|s|\geq b$. The sum is locally finite, periodic, odd,
and equals $s$ for $|s|\leq a$. In that region,
\[
\vartheta=\lambda\Theta\,\zeta\cdot x,\qquad
v=\frac{\Omega}{q}J\zeta\,(\zeta\cdot x),\qquad \varpi=\Omega.
\]
In particular
\begin{equation}
\nabla\vartheta(0,t)=\lambda\Theta\zeta,\qquad
Dv(0,t)=\frac{\Omega}{q}J\zeta\otimes\zeta.
\label{eq:bc-central-gradient}
\end{equation}
These two identities also hold for the sine profile since
$F(0)=0$ and $F'(0)=1$.

\subsection{Thermal diffusion and viscosity}

For the same general profile, subtracting the affine background from
the diffusive equations gives the exact scalar and vorticity residuals
\begin{align}
R_\theta&=
\left(\dot\Theta+\frac{\Omega}{\lambda q}J\zeta\cdot G\right)F(s)
-\kappa\lambda^2q\,\Theta F''(s),\label{eq:bc-general-diffusion-scalar}\\
R_\omega&=(\dot\Omega-\lambda\zeta_1\Theta)F'(s)
-\nu\lambda^2q\,\Omega F'''(s).
\label{eq:bc-general-diffusion-vorticity}
\end{align}
Indeed, $\Delta\vartheta=\lambda^2q\,\Theta F''(s)$ and
$\Delta\varpi=\lambda^2q\,\Omega F'''(s)$, while all advective
cancellations proved above still hold. Thus the general profile is
not automatically closed under the diffusive equations.

For the sine profile, $F''=-F$ and $F'''=-F'$, and the exact amplitude
equations are instead
\begin{equation}
\dot\Theta=-\kappa\lambda^2q\,\Theta
-\frac{J\zeta\cdot G}{\lambda q}\Omega,\qquad
\dot\Omega=\lambda\zeta_1\Theta-\nu\lambda^2q\,\Omega.
\label{eq:bc-diffusive-amplitudes}
\end{equation}
These equations also lift to the full diffusive momentum equation
without changing the prescribed background forces. In fact replace
$h$ in \eqref{eq:bc-pressure} by
\[
h_\nu=\dot C J\zeta+2CDJ\zeta+\nu\lambda^2q\,C J\zeta-\Theta e_2
\]
and put
$\pi_\nu=-(\zeta\cdot h_\nu)P(s)/(\lambda q)$.
The additional momentum term is exactly
$-\nu\Delta v=\nu\lambda^2q\,v$.
Moreover
\[
J\zeta\cdot h_\nu
=\frac{\dot\Omega+\nu\lambda^2q\,\Omega}{\lambda}
-\Theta\zeta_1=0
\]
by \eqref{eq:bc-diffusive-amplitudes}. The same orthogonal decomposition
and gradient calculation therefore give
$\nabla\pi_\nu=-h_\nu F(s)$, proving the asserted full momentum identity.

\subsection{The frozen matrix with every frequency factor retained}

Fix
\[
D=0,\qquad G=-Ae_2,\qquad A>0,\qquad
\zeta=r(\sin\varphi,\cos\varphi),\qquad r>0,\quad\varphi\in\mathbb R.
\]
Here $r,\lambda,A,\varphi$ are constants, and $k=\lambda r>0$
is the magnitude of the physical wavevector.
This is an exact stationary unforced background when
\[
u_b=0,\quad\theta_b=-Ax_2,\quad p_b=-\frac A2x_2^2:
\quad \nabla p_b=\theta_be_2,\quad f_{\theta,b}=f_{u,b}=0.
\]
Its affine temperature does not decay in space.

Since $J\zeta=r(-\cos\varphi,\sin\varphi)$, we have
$J\zeta\cdot G=-Ar\sin\varphi$ and $\zeta_1=r\sin\varphi$.
Equations \eqref{eq:bc-diffusive-amplitudes} become
\begin{equation}
\frac{d}{dt}\begin{pmatrix}\Theta\\\Omega\end{pmatrix}
=
B\begin{pmatrix}\Theta\\\Omega\end{pmatrix},\qquad
B=\begin{pmatrix}
-\kappa k^2&A\sin\varphi/k\\
k\sin\varphi&-\nu k^2
\end{pmatrix}.
\label{eq:bc-frozen-matrix}
\end{equation}
No wavelength or coefficient has been absorbed into a new unit system.

\begin{proposition}[Exact growth threshold for the frozen sine wave]
The two eigenvalues of \eqref{eq:bc-frozen-matrix} are
\begin{equation}
\mu_\pm=-\frac{(\kappa+\nu)k^2}{2}
\ \pm\
\sqrt{\frac{(\kappa-\nu)^2k^4}{4}+A\sin^2\varphi}.
\label{eq:bc-eigenvalues}
\end{equation}
There is a strictly positive eigenvalue if and only if
\begin{equation}
A\sin^2\varphi>\kappa\nu\,\lambda^4r^4.
\label{eq:bc-growth-threshold}
\end{equation}
\end{proposition}
\begin{proof}
Direct expansion yields
\[
\det(\mu\mathrm{Id}-B)
=(\mu+\kappa k^2)(\mu+\nu k^2)-A\sin^2\varphi
=\mu^2+(\kappa+\nu)k^2\mu+\kappa\nu k^4-A\sin^2\varphi.
\]
Completing the square gives \eqref{eq:bc-eigenvalues}.
The expression under the square root is nonnegative.
Since $\kappa,\nu\geq0$, $\mu_-\leq0$, and $\mu_+>0$ is equivalent
to the square root exceeding $(\kappa+\nu)k^2/2$.
Both quantities are nonnegative, so squaring is an equivalence.
Subtracting $(\kappa+\nu)^2k^4/4$ from the radicand gives
$A\sin^2\varphi-\kappa\nu k^4$. Finally $k^4=\lambda^4r^4$.
This proves both directions, including the strict inequality.
\end{proof}

For $\kappa=\nu=0$ and $0<\varphi<\pi$, the positive eigenvalue is
$\sqrt A\sin\varphi$ and its eigenline is
\[
\Omega=\frac{\lambda r}{\sqrt A}\Theta.
\]
Substitution into both rows of \eqref{eq:bc-frozen-matrix} proves
this statement directly. Thus along that line
\[
\Theta(t)=\Theta(0)e^{\sqrt A\sin\varphi\,t},\qquad
\Omega(t)=\frac{\lambda r}{\sqrt A}\Theta(0)
e^{\sqrt A\sin\varphi\,t}.
\]
Negative $\Theta(0)$ gives negative $\Omega$ and, by
\eqref{eq:bc-central-gradient}, a temperature-gradient increment
pointing along $-\zeta$. Its Euclidean magnitude is
$\lambda r|\Theta|$, whereas the sine temperature amplitude is
$|\Theta|$. Also
$Dv(0)=\Omega Je(\varphi)\otimes e(\varphi)$ for the original
$e(\varphi)=(\sin\varphi,\cos\varphi)$.

If $\kappa,\nu>0$ with $A$ fixed, the threshold fails whenever
$k^4\geq A/(\kappa\nu)$, because $\sin^2\varphi\leq1$.
If $\kappa=0$, $\nu>0$, and $\sin\varphi\ne0$, the positive rate is
still strictly positive, but the exact rationalized expression is
\[
\mu_+=
\frac{A\sin^2\varphi}
{\sqrt{\nu^2k^4/4+A\sin^2\varphi}+\nu k^2/2}
\leq \frac{A\sin^2\varphi}{\nu k^2}.
\]
The equality comes from multiplication by the conjugate, and the
bound uses a denominator at least $\nu k^2$.
Thus even in this one-diffusion case the inviscid growth rate cannot
be carried over unchanged.

A fixed finite-dimensional matrix in \eqref{eq:bc-frozen-matrix}
does not produce a finite-time singularity: its solution is
$\sum_{n=0}^{\infty}t^nB^n/n!$ times the initial vector, bounded by
$e^{|t|\|B\|}$ times that vector's norm on every bounded time interval.
The spatial fields constructed from it are smooth at every finite
time. A nonzero uncut wave also has no finite spatial $L^2$ norm:
on a strip of positive width where its periodic profile has
magnitude bounded below, the profile is constant in the unbounded
direction perpendicular to $\zeta$, so integrating its square in
that direction diverges. The affine background separately fails
the spatial decay requirements.

\subsection{Exact common damping through growth, steering and holding}

The released construction uses a growth interval followed by a rotation
that returns the new vorticity amplitude to zero. Its holding interval
then prepares the next layer. The following identity computes the effect
of equal diffusion coefficients on the sine-wave amplitude system for
the same prescribed affine background, including time-dependent phase
length and steering. It does not assume that the background generated
by a whole family of interacting layers is unchanged by viscosity.

Retain a compact interval $[t_0,t_1]$, the original $D,G,\zeta,\lambda$,
and $q=|\zeta|^2$. Define
\[
 A_0(t)=\begin{pmatrix}
 0&-(J\zeta\cdot G)/(\lambda q)\\
 \lambda\zeta_1&0
 \end{pmatrix},\qquad
 d(t,t_0)=\exp\left(-\nu\lambda^2\int_{t_0}^tq(s)\,ds\right).
\]
For $\kappa=\nu>0$, the exact viscous amplitude matrix is
$A_0(t)-\nu\lambda^2q(t)\mathrm{Id}$. Let $Z_E=(\Theta_E,\Omega_E)$
be the uniquely defined solution of $Z_E'=A_0Z_E$ with any real
initial pair. Then the viscous solution with that same pair is
\begin{equation}\label{eq:bc-common-damping}
 Z_\nu(t)=d(t,t_0)Z_E(t).
\end{equation}
Indeed $d'=-\nu\lambda^2qd$, and the ordinary product rule gives
\[
 (dZ_E)'=(A_0-\nu\lambda^2q\mathrm{Id})(dZ_E).
\]
Also $d(t_0,t_0)=1$. The existence and uniqueness argument for the
linear system given earlier in this section proves the assertion and
its converse after multiplication by $d^{-1}$. The scalar factor
commutes with $A_0(t)$ even when the matrices $A_0(t)$ at different
times do not commute with each other.

Since $d>0$ at each finite time, zeros and signs of both amplitude
components are preserved at exactly the same times in these two
systems. In particular, a return $\Omega_E(t_*)=0$ remains a return
$\Omega_\nu(t_*)=0$. During a subsequent interval where $\zeta_1=0$
and the initial vorticity is zero, uniqueness gives $\Omega_E=0$,
$\Theta_E$ constant, and
\[
 \Omega_\nu(t)=0,\qquad
 \Theta_\nu(t)=\Theta_\nu(t_*)
 \exp\left(-\nu\lambda^2\int_{t_*}^tq(s)\,ds\right).
\]
Thus the held vorticity remains zero while the held temperature
amplitude decays by this exact factor.

When $\Theta_E(t_0)\ne0$ and $\Theta_E(t_1)\ne0$, put
$g=|\Theta_E(t_1)|/|\Theta_E(t_0)|>0$.
The viscous amplitude gain on the complete chosen interval is exactly
\begin{equation}\label{eq:bc-cycle-gain}
 \frac{|\Theta_\nu(t_1)|}{|\Theta_\nu(t_0)|}
 =g\exp\left(-\nu\lambda^2\int_{t_0}^{t_1}q(s)\,ds\right).
\end{equation}
It exceeds one precisely when
$\log g>\nu\lambda^2\int_{t_0}^{t_1}q(s)\,ds$, by monotonicity
of the logarithm. In the frozen growing eigenline, with
$0<\varphi<\pi$, this gives the retained rate
$\sqrt A\sin\varphi-\nu\lambda^2r^2$ and exactly the equal-diffusion
case of \eqref{eq:bc-growth-threshold}.

For a general periodic phase profile, its Fourier mode $m$ instead
has the factor
\[
 d_m(t,t_0)=\exp\left(-\nu m^2\lambda^2
                                  \int_{t_0}^tq(s)\,ds\right).
\]
This follows by the same matrix product calculation after
$\partial_s^2e^{ims}=-m^2e^{ims}$. The sine factor $d_1$ applies to
modes $|m|=1$. Nonzero modes with distinct absolute frequencies have
different damping factors for $t>t_0$, so their sum generally has no
common scalar damping factor;
the full Fourier evolution is the one constructed in
Section~\ref{sec:localized-layer}.

That section also gives exact finite localized corrections with their
remaining residual. Section~\ref{sec:swirl-carrier} gives the full
three-dimensional swirl transfer and its computed force. Neither the
frozen eigenvalue nor the common-damping identity establishes a
singular limit of the coupled, localized sequence.
